\section{功能需求}

本章给出系统的核心功能需求。所有功能项均采用“需求编号 + 功能名称 + 描述 + 前置条件 + 预期结果”的形式表述，默认视为本版本范围内的必选需求。

\subsection{账号与身份认证}
本模块负责用户生命周期的起始环节，为后续好友关系、社区管理和实时交互提供统一的身份基础。

\Requirement{FR-01}{用户注册}{
系统应支持未登录用户通过邮箱或手机号、用户名和密码注册平台账号，并建立基础用户档案。
}{
用户尚未登录，且提交的邮箱或手机号未被平台占用。
}{
系统创建新账号，记录注册时间与基础状态，向用户返回注册结果，并在需要时触发身份验证流程。
}

\Requirement{FR-02}{登录与身份验证}{
系统应支持注册用户通过合法凭证登录，并校验账号状态、密码正确性和必要的身份验证条件。
}{
用户已完成注册且账号处于可登录状态。
}{
系统创建有效会话，加载用户基础资料、好友状态和通知摘要；若登录失败，应返回明确失败原因。
}

\Requirement{FR-03}{密码找回与重置}{
系统应支持用户通过受信任的验证方式找回密码，并在验证通过后完成密码重置。
}{
用户提供与账号绑定的验证信息，并通过平台验证。
}{
系统允许用户设置新密码，使旧密码失效，并保留必要的安全日志记录。
}

\subsection{个人资料、在线状态与社交关系}
本模块描述用户在平台层面的个人信息维护、状态展示和一对一社交行为。

\Requirement{FR-04}{个人资料维护}{
系统应支持用户查看和修改昵称、头像、个性签名等基础资料，并允许保存个性化展示信息。
}{
用户已成功登录。
}{
系统保存资料变更内容，并在需要展示用户资料的界面中同步更新。
}

\Requirement{FR-05}{在线状态同步}{
系统应支持展示用户在线、离开、忙碌、隐身和离线等状态，并在好友列表、私聊窗口及社区成员列表中实时同步。
}{
用户已登录，且客户端与服务端保持有效连接。
}{
相关用户可在权限允许的前提下看到最新状态变化，状态变化事件被实时分发。
}

\Requirement{FR-06}{好友申请与处理}{
系统应支持用户搜索或发起好友申请，并允许接收方接受、拒绝或忽略申请。
}{
双方均为平台注册用户，且尚未建立好友关系。
}{
系统更新好友申请状态；当申请被接受时，双方建立好友关系，并可进入私聊能力范围。
}

\Requirement{FR-07}{私聊会话}{
系统应支持已建立好友关系的用户发起一对一私聊会话，并保留会话历史记录和未读状态。
}{
双方已经互为好友，且至少一方主动发起私聊。
}{
系统创建或复用私聊会话，允许双方发送消息，并更新最近会话时间与未读计数。
}

\subsection{社区与成员管理}
本模块定义围绕社区空间的创建、加入、退出和成员层面的基础管理需求。

\Requirement{FR-08}{创建社区}{
系统应支持已登录用户创建社区，并设置社区名称、图标、简介和默认权限结构。
}{
用户已登录，且满足社区创建条件。
}{
系统生成新的社区实体，并将创建者设为社区所有者，同时创建默认频道或默认角色配置。
}

\Requirement{FR-09}{加入、邀请与退出社区}{
系统应支持用户通过邀请码、邀请链接或搜索入口加入社区，并允许成员主动退出社区。
}{
用户已登录，邀请码或加入入口有效，且未被社区限制访问。
}{
系统更新成员关系，加入成功后用户获得默认成员角色；退出成功后移除其社区访问权限。
}

\Requirement{FR-10}{社区成员列表与基础管理}{
系统应支持社区内成员列表展示，并允许具备权限的管理者查看成员信息、移除成员或调整其基础状态。
}{
用户已加入社区，且对目标操作拥有相应管理权限。
}{
系统根据操作结果更新成员清单、成员状态和可见权限，并将变更同步到相关客户端。
}

\subsection{频道管理}
本模块定义社区内部文本频道和语音频道的组织方式及访问控制要求。

\Requirement{FR-11}{频道创建与配置}{
系统应支持具备权限的用户创建文本频道或语音频道，并配置名称、主题、排序和可见性。
}{
用户已加入目标社区，且拥有频道管理权限。
}{
系统创建频道记录，更新社区频道列表，并按配置展示在客户端导航区域。
}

\Requirement{FR-12}{频道访问控制}{
系统应支持按角色或成员维度限制频道的查看、发言、接入语音等操作权限。
}{
目标社区已存在角色或成员权限规则，且管理者具备权限配置能力。
}{
系统在用户访问频道或执行相关操作时进行权限校验，仅向有权用户展示频道内容或允许其互动。
}

\subsection{文本消息}
本模块聚焦文本频道和私聊中的消息生产、分发、历史查看及内容治理要求。

\Requirement{FR-13}{文本消息发送与附件上传}{
系统应支持用户在文本频道或私聊中发送文字消息、图片或文件附件，并支持消息内提及其他用户。
}{
用户已登录并拥有当前会话的发送权限，附件类型与大小满足平台限制。
}{
系统保存消息内容和附件元数据，将消息实时分发给会话参与者或频道成员，并对被提及用户生成提醒。
}

\Requirement{FR-14}{历史消息加载与未读定位}{
系统应支持按时间顺序加载历史消息，并在用户重新进入会话时定位到最近未读位置。
}{
目标文本频道或私聊会话存在历史消息记录。
}{
系统按分页或滚动方式返回消息历史，并更新用户在当前会话中的最近读取位置。
}

\Requirement{FR-15}{消息撤回与删除}{
系统应支持消息发送者在规定条件下撤回自身消息，并允许具备管理权限的用户删除违规消息。
}{
目标消息存在，且操作人满足消息拥有权或频道管理权限要求。
}{
系统更新消息可见状态，向相关客户端同步删除或撤回结果，并保留必要的操作记录。
}

\subsection{语音频道互动}
本模块用于定义多人语音加入、退出及语音状态同步需求。

\Requirement{FR-16}{加入与退出语音频道}{
系统应支持拥有权限的用户加入语音频道、保持在线连接，并在结束会话后主动退出。
}{
用户已登录且有权访问目标语音频道。
}{
系统建立或释放语音会话状态，更新频道成员在线列表，并同步用户当前所处的语音频道信息。
}

\Requirement{FR-17}{语音状态同步}{
系统应支持用户在语音频道中执行静音、取消静音、关闭麦克风、切换设备状态等操作，并向频道成员同步状态变化。
}{
用户已处于某个语音频道的有效会话中。
}{
系统实时更新成员语音状态，使频道内其他成员看到一致的状态展示结果。
}

\subsection{通知与提醒}
本模块用于保障用户能够及时获知与其相关的重要交互事件。

\Requirement{FR-18}{通知生成与已读管理}{
系统应针对好友申请、私聊新消息、社区消息提及、频道未读和管理动作结果生成通知，并支持通知已读处理。
}{
系统发生与用户相关的业务事件，且用户具备接收该事件通知的条件。
}{
系统为目标用户生成通知记录，更新未读计数，并在用户查看后将通知状态更新为已读。
}

\subsection{角色与权限控制}
本模块用于定义社区内的角色能力边界和具体权限校验行为。

\Requirement{FR-19}{角色定义与分配}{
系统应支持社区所有者或管理员创建角色，并为角色配置名称、优先级、颜色和权限集合，再将角色分配给成员。
}{
用户处于目标社区内，并拥有角色管理权限。
}{
系统保存角色配置与成员角色映射关系，使角色设置立即作用于频道访问和管理操作。
}

\Requirement{FR-20}{权限校验与拦截}{
系统应在用户执行查看频道、发送消息、加入语音、删除消息、管理成员等操作前进行权限校验。
}{
用户发起需要权限判断的业务请求。
}{
当权限满足时系统允许继续执行；当权限不足时系统拒绝请求并返回明确提示，不得造成越权访问。
}
